\documentclass[12pt,a4paper]{article}
\usepackage[utf8]{inputenc}
\usepackage[vietnamese]{babel}
\usepackage[T5]{fontenc}
\usepackage{amsmath,amssymb,amsfonts}
\usepackage{fancyhdr}
\usepackage{geometry}
\usepackage{tcolorbox}

\geometry{
	a4paper,
	left=20mm,
	right=20mm,
	top=20mm,
	bottom=20mm
}

\pagestyle{fancy}
\fancyhf{}
\lhead{\small\textit{Bổ đề Hình học: Phân giác, Chùm điều hòa \& Đường đối trung}}
\rhead{\small\textbf{Integra Typesetter}}
\rfoot{\small Trang \thepage}
\renewcommand{\headrulewidth}{0.4pt}
\renewcommand{\footrulewidth}{0.4pt}

\begin{document}

\begin{center}
	{\Large \textbf{CÁC BỔ ĐỀ HÌNH HỌC PHẲNG CHUYÊN SÂU}}\\[2mm]
	{\small \textit{(Đường phân giác -- Chùm điều hòa trực giao -- Điểm tiếp xúc đường tròn nội tiếp)}}
\end{center}
\vspace{4mm}

\begin{tcolorbox}[colback=blue!5!white,colframe=blue!75!black,title=\textbf{Bổ đề 1: Tính chất Giao điểm của Phân giác và Đường trung bình}]
Cho tam giác $ABC$ có đường tròn nội tiếp $(I)$ tiếp xúc với ba cạnh tại $J, E, F$. Phân giác trong $BI, CI$ lần lượt cắt đường trung bình $MN$ ($M \in AB, N \in AC$) tại $K, H$. Khi đó:
\begin{enumerate}
    \item $K \in JE$ và $AK \perp BI$.
    \item $H \in JF$ và $AH \perp CI$.
    \item Tứ giác $AKJH$ là hình bình hành, suy ra $JA$ đi qua trung điểm của $HK$.
\end{enumerate}
\end{tcolorbox}

\vspace{3mm}

\begin{tcolorbox}[colback=teal!5!white,colframe=teal!75!black,title=\textbf{Bổ đề 2: Chùm điều hòa Trực giao}]
Cho hai chùm điều hòa bốn tia $O_1(ABCD) = -1$ và $O_2(A'B'C'D') = -1$. Nếu ba cặp tia tương ứng vuông góc:
\[
O_1A \perp O_2A', \quad O_1B \perp O_2B', \quad O_1C \perp O_2C'
\]
thì cặp tia thứ tư cũng vuông góc: $O_1D \perp O_2D'$.
\end{tcolorbox}

\vspace{3mm}

\begin{tcolorbox}[colback=orange!5!white,colframe=orange!75!black,title=\textbf{Bổ đề 3: Sự đồng quy giữa Trung tuyến và Tiếp điểm}]
Cho tam giác $ABC$ nội tiếp đường tròn có $(I)$ tiếp xúc $BC, CA, AB$ tại $D, E, F$. $K$ là trung điểm $BC$. Khi đó ba đường thẳng $AK, ID, EF$ đồng quy.
\end{tcolorbox}

\end{document}
